\section{The whole procedure on a small example}\label{sec:matching-trace}
Let three students $a$, $b$, $c$ accept tasks from $1$, $2$, $3$ as follows:
\[
N(a)=\{1,2\},\qquad N(b)=\{1\},\qquad N(c)=\{2,3\}.
\]
We run the procedure of Section~\ref{ch08:sec:search}, taking the unmatched students as roots in the order $a$, $c$, $b$, and inspecting the edges at each student in increasing order of task number.
\begin{enumerate}
\item Starting from the empty matching, the search from $a$ finds the unused task~$1$ at once. The augmenting path is the single edge $a$--$1$, and flipping gives $M=\{a\text{--}1\}$.
\item Likewise, the search from $c$ finds the unused task~$2$ at once, and flipping gives $M=\{a\text{--}1,\ c\text{--}2\}$.
\item Now $b$ is unmatched, and its only task is held by $a$. The first two rounds made exactly the choices of a greedy procedure, which would now be stuck. The search from $b$ reaches the following vertices, in this order.
\end{enumerate}
\begin{center}
\begin{tabularx}{\linewidth}{@{}llX@{}}
\toprule
Vertex reached & Predecessor & Reason\\
\midrule
$b$ & --- & the root of the search\\
task $1$ & $b$ & inspecting $b$: the edge $b$--$1$ is unmatched\\
$a$ & task $1$ & task $1$ is held by $a$\\
task $2$ & $a$ & inspecting $a$: the edge $a$--$2$ is unmatched ($a$--$1$ is matched)\\
$c$ & task $2$ & task $2$ is held by $c$\\
task $3$ & $c$ & inspecting $c$: the edge $c$--$3$ is unmatched ($c$--$2$ is matched)\\
\bottomrule
\end{tabularx}
\end{center}
Task~$3$ is unused, so the search succeeds. Following the predecessors backward from task~$3$ gives $3$, $c$, $2$, $a$, $1$, $b$, and read forward this is the augmenting path
\[
b\;\text{--}\;1\;\text{--}\;a\;\text{--}\;2\;\text{--}\;c\;\text{--}\;3.
\]
The picture shows the flip.

\begin{center}
\begin{tikzpicture}[x=0.88in,y=0.5in,every node/.style={font=\small}]
\foreach \x/\name in {0/b,1/1,2/a,3/2,4/c,5/3}
 {\node[circle,draw=ink,minimum size=7mm,inner sep=1pt] (v\x) at (\x,0) {$\name$};}
\draw[accent,dashed,thick] (v0)--(v1) (v2)--(v3) (v4)--(v5);
\draw[ink,line width=1.2pt] (v1)--(v2) (v3)--(v4);
\node[align=center] at (2.5,-.8) {Solid: matched edges, removed by the flip.\\Dashed: unmatched edges, added by the flip.};
\end{tikzpicture}
\end{center}

After the flip, the matching $\{b\text{--}1,\ a\text{--}2,\ c\text{--}3\}$ has three edges. Every student is matched, so the procedure stops and reports this assignment. Anyone can check it: $1\in N(b)$, $2\in N(a)$, $3\in N(c)$, and the three tasks are different.

Now remove task~$3$ from $c$'s options, so that $N(c)=\{2\}$, and run the procedure again in the same order. The first two rounds again give $M=\{a\text{--}1,\ c\text{--}2\}$, and the search from $b$ again reaches task~$1$, $a$, task~$2$, and $c$. But now the only edge at $c$ is the matched edge $c$--$2$, so inspecting $c$ reaches nothing new, and the search fails with
\[
A=\{b,a,c\},\qquad B=\{1,2\}.
\]
As Proposition~\ref{ch08:prop:shortage} predicts, $|B|=|A|-1$ and
\[
N(A)=N(a)\cup N(b)\cup N(c)=\{1,2\}\cup\{1\}\cup\{2\}=\{1,2\}=B.
\]
So $|N(A)|=2<3=|A|$, and $A$ is a shortage set: no order of choices could have succeeded.
